%! Author = eniko
%! Date = 2025. 10. 11.

% Preamble
\documentclass[11pt]{article}

% Packages
\usepackage{amsmath}

% Document
\begin{document}
:sectnums:
:sectnumlevels: 3
:stem: latexmath

== Geometriai Brown mozgás

=== Bevezetés

A tudományos modellezés kezdetben a klasszikus fizika által inspirált determinizmusra épült.
Ennek alapvető matematikai eszközei a közönséges differenciálegyenletek (ODE-k), amelyek feltételezik, hogy egy rendszer jövőbeli állapota a jelenlegi állapotból egyértelműen kiszámítható. Egy ODE adott kezdőfeltételek mellett egyetlen, előre meghatározott trajektóriát ír le.
Klasszikus példák erre a bolygók mozgása a Nap körül vagy egy ideális, súrlódásmentes rugó által mozgatott test pályája, ahol a jövőbeli pozíció és sebesség pontosan kiszámítható.
Azonban a komplex biológiai, pénzügyi vagy társadalmi rendszerek elemzése során nyilvánvalóvá vált, hogy valós rendszerben a bizonytalanság és a véletlenszerűség nem elhanyagolható, hanem a dinamika szerves és megkerülhetetlen része.
Egy részvényárfolyam, egy járvány terjedése vagy a bűnözési ráta alakulása olyan folyamatok, amelyeket részben vagy teljesen véletlen hatások befolyásolnak.
Ezen jelenségek leírására született meg a sztochasztikus (vagy véletlenszerű) folyamat fogalma, amely olyan időben végbemenő folyamatot jelöl, amelyet valószínűségi változók jellemeznek.
Az SDE-k a sztochasztikus folyamatok egy fontos osztályának dinamikáját (időbeli fejlődését) leíró matematikai eszközök, amelyek lehetővé teszik a véletlenszerű komponensek beépítését a dinamikai modellekbe.
Míg egy ODE egy adott kezdőpontból egyetlen, jól definiált megoldást jelöl ki, addig egy SDE egyetlen kezdőpontból a lehetséges jövőbeli kimenetelek végtelen sokaságát vázolja fel.
Ellentétben a közönséges differenciálegyenletek (ODE-k) egyedi, determinisztikus trajektóriát adó megoldásaival, egy sztochasztikus differenciálegyenlet (SDE) megoldása maga egy sztochasztikus folyamat.
Ez a folyamat a lehetséges trajektóriák halmazán definiál egy valószínűségi mértéket.
Tehát a megoldás nem egyetlen függvény, hanem a trajektóriák teljes sokasága, valamint az ezen sokaság felett értelmezett valószínűségi eloszlás, amely meghatározza az egyes útvonal-családok bekövetkezésének valószínűségét.

=== A Sztochasztikus Folyamatok Elméleti Alapjai

==== Definíciók és alapfogalmak

.Valószínűségi tér

Egy valószínűségi tér stem:[(\Omega, \mathcal{F}, P)] hármasa tartalmazza:

* stem:[\Omega] *(Eseménytér)*: A lehetséges kimenetelek halmaza.
* stem:[\mathcal{F}] *(stem:[\sigma]-algebra)*: stem:[\Omega] részhalmazainak olyan rendszere, amely tartalmazza az stem:[\emptyset] üres halmazt; komplementer-zárt; és megszámlálható unióra zárt. Az stem:[(\Omega, \mathcal{F})] párost mérhető térnek nevezzük.
* stem:[P] *(Valószínűségi Mérték)*: Egy stem:[P: \mathcal{F} \to \mathbb{R}] függvény, amely kielégíti a stem:[P(\emptyset)=0] és stem:[P(\Omega)=1] feltételeket, és megszámlálható additivitással rendelkezik diszjunkt halmazok esetén.

*Definíció*: Borel σ-algebra: A R nyitott részhalmazai generálják a Borel σ-algebrát (stem:[\mathcal{B}]).

*Definíció*: A függvények egy halmaza által *generált σ-algebra* a legkisebb σ-algebra, amelyre vonatkozóan ezen függvények mind mérhetők.

.Abszolút folytonosság

Tegyük fel, hogy adott két valószínűségi mérték stem:[\mathcal{P}],stem:[\mathcal{Q}]: stem:[\mathcal{F}] -> [0,1] ugyanazon mérhető téren stem:[(\Omega, \mathcal{F})].

*Definíció*: Azt mondjuk, hogy a stem:[Q] mérték abszolút folytonos a stem:[P]-re nézve, ha létezik egy integrálható véletlen változó stem:[f] stem:[\Omega \to \mathbb{R}] melyre minden A stem:[\in \mathcal{F}] esetén :

\[
\mathcal{Q}(A) = \int_A f (\omega) \, d\mathcal{P}(\omega)
\]

Ekkor stem:[f] a stem:[\mathcal{Q}] mérték stem:[\mathcal{P}]-re vonatkozó sűrűségfüggvénye. Ebben az esetben ha stem:[E^{\mathcal{P}}] és stem:[E^{\mathcal{Q}}] a két mérték várható értékei stem:[\mathcal{P}] és stem:[\mathcal{Q}] szerint, akkor minden stem:[\mathcal{Q}]- szerint integrálható véletlen változóra stem:[X]:
\[
E^{\mathcal{Q}}[X] = \int_{\Omega} Xd\mathcal{Q} = \int_{\Omega} X f d\mathcal{P} = E^{\mathcal{P}}[X f]
\]

.Eloszlás és sűrűségfüggvény

Egy valószínűségi változó stem:[X: \Omega \to \mathbb{R}] eloszlásfüggvénye stem:[F_X: \mathbb{R} \to [0,1]] a következőképpen van definiálva:

\[
F_X(x) = \mathcal{P}(\{\omega \in \Omega : X(\omega) \le x\})
\]

*Definició* : Egy valószinűségi változó stem:[X] folytonos, ha a valószínűségi mérték stem:[\mathcal{P}_X] abszolút folytonos a Lebesgue mértékre nézve.

Ebben az esetben az eloszlásfüggvénye differenciálható:

\[
F_X^{\prime}(x) = f_X(x)
\]

*Definíció* : Ekkor stem:[f_X: \mathbb{R} \to [0, \infty)] a valószínűségi változó stem:[X] sűrűségfüggvénye. A sűrűségfüggvény kielégíti a következő feltételeket:

1. *Nemnegativitás*: Minden x ∈ ℝ esetén f_X(x) ≥ 0.
2. *Integrálhatóság*: A sűrűségfüggvény integráljára teljesül, hogy

\[
\int_{-\infty}^{\infty} f_X(x) \, dx = 1
\]

Ha stem:[A \in \mathcal{B}(\mathbb{R})], akkor annak a valószínűsége, hogy stem:[X] értéke stem:[A]-ba esik az alábbi integrállal számítható ki:

\[
\mathcal{P}(X \in A) = \int_A f_X(x) \, dx
\]

.Feltételes várható érték
Legyen stem:[X: \Omega \to \mathbb{R}] egy integrálható valószínűségi változó a (stem:[\Omega, \mathcal{F}, \mathcal{P}]) valószinűségi téren.
A stem:[\mathcal{G} \subseteq \mathcal{F}] egy σ-algebra.

*Definíció*: A stem:[X] valószínűségi változó stem:[\mathcal{G}]-re vonatkozó feltételes várható értéke egy stem:[\mathcal{G}]-mérhető valószínűségi változó stem:[E[X|\mathcal{G}\]: \Omega \to \mathbb{R}]
amely kielégíti a következő feltételt: minden korlátos stem:[\mathcal{G}]-mérhető valószínűségi változóra stem:[Y]:
\[
E[XY] = E[E[X|\mathcal{G}]Y]
\]

Ha stem:[Y = 1_B] egy indikátor függvény egy stem:[B \in \mathcal{G}] eseményre, akkor minden stem:[B \in \mathcal{G}] teljesül:
\[
\int _B E[X|\mathcal{G}] d\mathcal{P} = \int_B X d\mathcal{P}
\]


=== Sztochasztikus folyamatok:

Tekintsünk egy olyan, időben „véletlenszerűen” változó, valós értékű mennyiségre, mint amilyen például a szél sebességének egy komponense egy meteorológiai állomáson,
vagy egy Brown-mozgást végző pollenszemcse hely-, illetve sebességkoordinátája.
Ezeket a folyamatokat egy stem:[ X: (0, \infty) \times \Omega \to \mathbb{R} ] mérhető függvénnyel írhatjuk le, ahol stem:[\Omega] egy valószínűségi tér
és stem:[X]-et sztochasztikus folyamatnak nevezzük. stem:[X(t,ω)] a folyamat értéke a t időpontban, a stem:[\omega \in \Omega] kimenetel esetén.
 A folyamat általában a stem:[t \in (0, \infty)] időintervallumon (vagy más időintervallumon, pl. diszkrét időpontokon) van definiálva.
Egy sztochasztikus folyamatra kétféleképpen is gondolhatunk. Először is, rögzítve egy stem:[\omega \in \Omega] elemet, egy időfüggvényt kapunk:

\[
X_{\omega} : t \mapsto X(t, \omega),
\]

amelyet a folyamat trajektóriájának nevezünk. Ebből a szempontból a folyamat az stem:[\{X_{\omega} : \omega \in \Omega\}] függvények gyűjteménye, a valószínűségi mérték pedig a trajektóriák terén értelmezett mérték.

Másrészről, rögzítve egy stem:[t \in [0, \infty)] időpontot, egy, az stem:[\Omega] valószínűségi téren értelmezett valószínűségi változót kapunk:

\[
X_t : \omega \mapsto X(t, \omega).
\]

Ebből a nézőpontból a folyamat a stem:[t] időváltozó által indexelt stem:[\{X_t : 0 \le t < \infty\}] valószínűségi változók gyűjteménye.
A valószínűségi mérték ezen valószínűségi változók együttes eloszlását  írja le.

.Markov folyamatok

*Definíció*: Egy sztochasztikus folyamat stem:[X = \{X_t : t \ge 0\}] Markov-folyamat, ha minden stem:[0 \le s < t] és minden Borel-mérhető stem:[f : \mathbb{R} \to \mathbb{R}] függvény esetén, ahol stem:[E|f(X_t)| < \infty], teljesül a következő feltétel:
\[E[f(X_t) | \mathcal{F}_s] = E[f(X_t) | X_s]\]
Ahol stem:[\mathcal{F}_s] a stem:[X]-folyamat stem:[s] időpontig ismert információit tartalmazó σ-algebra.
Tehát a Markov-folyamat jövőbeli állapota csak a jelenlegi állapottól függ, és nem függ a múltbeli állapotoktól.

Egy X folyamat Markov-folyamat, ha bármely stem:[0 \le t_1 \lt t_2 \lt \dots \lt t_{n+1}] időpont esetén a feltételes valószínűségi sűrűségfüggvény csak a legutolsó időponttól függ, azaz:

\[
p(x_{n+1}, t_{n+1} \mid x_n, t_n; \dots; x_2, t_2; x_1, t_1) = p(x_{n+1}, t_{n+1} \mid x_n, t_n)
\]

Ebből következik, hogy az együttes valószínűségi sűrűség feltételes formája felírható átmeneti sűrűségek szorzataként:
\[
p(x_n, t_n; \dots; x_2, t_2 \mid x_1, t_1) = p(x_n, t_n \mid x_{n-1}, t_{n-1}) \cdots p(x_2, t_2 \mid x_1, t_1).
\]

Így egy folytonos Markov-folyamat összes együttes véges dimenziós valószínűségi sűrűségfüggvényét meg lehet határozni az átmeneti sűrűségfüggvény stem:[p(x,t \mid y,s)] és a kezdeti érték stem:[X_0] sűrűségfüggvénye stem:[p_0(x)] segítségével.
Például az stem:[X_t] egypontos sűrűsége a következőképpen adható meg:

\[
p(x,t) = \int_{-\infty}^{\infty} p(x,t \mid y,0) p_0(y) dy.
\]

Egy Markov-folyamat átmeneti valószínűségei nem tetszőlegesek, hanem kielégítik a Chapman-Kolmogorov egyenletet. Folytonos Markov-folyamat esetén ez az egyenlet a következő:
\[
p(x,t \mid y,s) = \int_{-\infty}^{\infty} p(x,t \mid z,r) p(z,r \mid y,s) dz
\]
bármely stem:[s < r < t] esetén.

Jelentése: Ahhoz, hogy a folyamat az stem:[s] időpontban lévő stem:[y] pontból a stem:[t] időpontban lévő stem:[x] pontba jusson, át kell haladnia egy tetszőleges köztes stem:[r] időpontban valamilyen stem:[z] ponton.

Egy folytonos Markov-folyamat időben homogén (time-homogeneous), ha az átmeneti sűrűsége csak az időbeli különbségektől függ:
\[
p(x,t \mid y,s) = p(x, t-s \mid y,0)
\]

.Brown-mozgás

A standard, egydimenziós Brown-mozgás (Wiener-folyamat), amelyet gyakran stem:[W_t] vagy stem:[B_t] szimbólummal jelölnek, egy olyan időben folytonos sztochasztikus folyamat, amely az alábbi négy tulajdonságnak tesz eleget:

* A folyamat a stem:[t=0] időpontban a nullából indul, 1 valószínűséggel, azaz stem:[B_0 = 0].
* A folyamat pályái, azaz a stem:[t \mapsto W_t(\omega)] leképezések minden egyes stem:[\omega] elemi eseményre, majdnem biztosan (azaz 1 valószínűséggel) folytonosak. Ez egy kulcsfontosságú tulajdonság, amely megkülönbözteti a Wiener-folyamatot a diszkrét idejű véletlen bolyongásoktól.
* Bármely stem:[0 \le s_1 < t_1 \le s_2 < t_2] időpontsorozatra a stem:[B_{t_1} - B_{s_1}] és stem:[B_{t_2} - B_{s_2}] valószínűségi változók függetlenek egymástól. Ez a tulajdonság azt jelenti, hogy a folyamat jövőbeli változásai függetlenek a múltbeli változásoktól.
* Bármely stem:[s < t] időpontpárra a stem:[B_t - B_s] növekmény normális eloszlású, 0 várható értékkel és stem:[t-s] szórásnégyzettel. Formálisan: stem:[B_t - B_s \sim \mathcal{N}(0, t-s)]. Ez a feltétel magában foglalja, hogy a növekmények statisztikai tulajdonságai csak az időkülönbség hosszától függenek (stacionaritás), nem pedig annak abszolút helyétől az időtengelyen.











Egy közönséges differenciálegyenlet az alábbi formában írható fel:
[latexmath]
++++
\frac{dx(t)}{dt} = f(t, x(t))
++++



== Példa sztochasztikus differenciálegyenletre
Ha megengedünk némi bizonytalanságot egy differenciálegyenlet együtthatóiban, akkor egy realisztikusabb modelt tudunk felépíteni.
Tekintsünk egy egyszerű populációs modellt, ahol a populáció növekedési rátája arányos a populáció méretével.
Ezt a következő differenciálegyenlettel írhatjuk le:

[latexmath]
++++
\frac{dN(t)}{dt} = a(t)N(t)
++++

Ahol stem:[N(t)] a populáció mérete időben, és stem:[a(t)] a növekedési ráta az idő függvényében.

Ha a növekedési rátát nem ismerjük teljesen pontosan, akkor helyettesíthetjük egy determinisztikus és egy sztochasztikus komponens összegével:

[latexmath]
++++
a(t) = r(t) + n(t)
++++

Ahol stem:[r(t)] a determinisztikus komponens, stem:[n(t)] pontos viselkedését nem ismerjük, csak a valószínűségi eloszlását.



A standard, egydimenziós Wiener-folyamat, amelyet gyakran latexmath:[W_t] vagy latexmath:[B_t] szimbólummal jelölnek, egy olyan időben folytonos sztochasztikus folyamat, amely az alábbi négy tulajdonságnak tesz eleget:

* *Kiindulópont:* A folyamat a latexmath:[t=0] időpontban a nullából indul, 1 valószínűséggel. Matematikailag: latexmath:[W_0 = 0].
* *Folytonos trajektóriák:* A folyamat pályái, azaz a latexmath:[t \mapsto W_t(\omega)] leképezések minden egyes latexmath:[\omega] elemi eseményre, majdnem biztosan (azaz 1 valószínűséggel) folytonosak. Ez egy kulcsfontosságú tulajdonság, amely megkülönbözteti a Wiener-folyamatot a diszkrét idejű véletlen bolyongásoktól.
* *Független növekmények:* Bármely latexmath:[0 \le s_1 < t_1 \le s_2 < t_2] időpontsorozatra a latexmath:[W_{t_1} - W_{s_1}] és latexmath:[W_{t_2} - W_{s_2}] valószínűségi változók függetlenek egymástól. Ez a tulajdonság azt jelenti, hogy a folyamat jövőbeli változásai függetlenek a múltbeli változásoktól; a folyamat "nem emlékszik" a múltjára. Ez a Markov-tulajdonság egy erősebb formája.
* *Stacionárius, normális eloszlású növekmények:* Bármely latexmath:[s < t] időpontpárra a latexmath:[W_t - W_s] növekmény normális eloszlású, 0 várható értékkel és latexmath:[t-s] szórásnégyzettel. Formálisan: latexmath:[W_t - W_s \sim \mathcal{N}(0, t-s)]. Ez a feltétel magában foglalja, hogy a növekmények statisztikai tulajdonságai csak az időkülönbség hosszától függenek (stacionaritás), nem pedig annak abszolút helyétől az időtengelyen.

==== Brown-mozgás

Egy általános, Itô-típusú SDE a következő szimbolikus formában írható fel:
[latexmath]
++++
dX(t) = \mu(t, X(t))dt + \sigma(t, X(t))dB(t)
++++
Ez az egyenlet valójában egy integrálegyenlet rövidítése, és két fő részből áll:

* *Drift tag*: a folyamat változásának determinisztikus, előre jelezhető komponense egy infinitezimális stem:[dt] időlépés alatt.
* *Diffúziós tag* (stem:[\sigma(X_t, t)dW_t]): Ez a tag a sztochasztikus, előre nem látható komponenst képviseli. A hatásának nagyságát a stem:[\sigma(X_t, t)] volatilitás- vagy diffúziós függvény skálázza, amely függhet a rendszer aktuális állapotától és az időtől is. Ez a tag visz véletlenszerűséget a rendszer leírásába

A fent leírt folyamatot Brown-mozgásnak nevezzük, ha stem:[\mu(t, X(t))] és stem:[\sigma(t, X(t))] konstansok.


[latexmath]
++++
dX_t = \mu dt + \sigma dW_t
++++


==== Itô integrál

==== Itô formula

==== Geometriai Brown-mozgás

A pénzügyi modellezésben és más területeken, ahol a vizsgált mennyiségek (pl. árak, populációk mérete) nem vehetnek fel negatív értéket és a növekedésük arányos a jelenlegi méretükkel, a leggyakrabban használt modell a geometriai Brown-mozgás (GBM).
Az SDE, amely ezt a folyamatot leíró SDE a következő :
[latexmath]
++++
dS(t) = \mu S(t) dt + \sigma S(t) dB(t)
++++
Ahol stem:[S(t)] a vizsgált mennyiség (pl. árfolyam), stem:[\mu] a drift (várható növekedési ráta), stem:[\sigma] a volatilitás (a véletlenszerű ingadozás mértéke), és stem:[B(t)] egy standard Brown-mozgás.
A geometriai Brown-mozgás egyenletének egyik nagy előnye, hogy létezik rá explicit, zárt alakú megoldás.
Ez a megoldás a következőképpen néz ki:
[latexmath]
++++
S(t) = S(0) \exp\left( \left(\mu - \frac{\sigma^2}{2}\right)t + \sigma B(t) \right)
++++
Ahol stem:[S(0)] a kezdeti érték.
Ez a megoldás azt mutatja, hogy a folyamat logaritmusa normális eloszlású, ami azt jelenti, hogy maga a folyamat lognormális eloszlású.
A megoldás levezetése az Itô formula alkalmazásával történik.




=== Alkalmazás

A geometriai Brown-mozgás (GBM) modelljét egy valós adatsoron, a bűnözési statisztikákon alkalmaztuk. A cél a napi bűnesetek számának előrejelzése volt egy kiválasztott kerületben a historikus adatok alapján. Az elemzés a `Geometriai Brown mozgas.ipynb` Jupyter notebookban található.

==== A modell felépítése és diszkretizálása

Az előrejelzéshez a GBM zárt alakú megoldását használtuk, amely két fő komponensből áll:

*   **Drift:** A hosszabb távú trendet írja le. A modellben a driftet a stem:[(\mu - \frac{1}{2}\sigma^2)t] tag képviseli, amely a folyamat determinisztikus növekedését vagy csökkenését határozza meg.
*   **Diffúzió:** A rövidebb távú, véletlenszerű ingadozásokat modellezi. Ezt a stem:[\sigma B(t)] tag fejezi ki, ahol a stem:[B(t)] Brown-mozgás a véletlen sokkok forrása, stem:[\sigma] volatilitás pedig skálázza ezeket a sokkokat.

A szimulációhoz a folytonos idejű megoldást diszkrét időlépésekre bontottuk. Az stem:[S_k] értéket az stem:[S_{k-1}] értékből származtatjuk:

[latexmath]
++++
S_k = S_{k-1} \cdot \exp\left( \left(\mu - \frac{\sigma^2}{2}\right)\Delta t + \sigma \sqrt{\Delta t} Z_k \right)
++++

ahol stem:[Z_k] egy standard normális eloszlású véletlen változó. A teljes idősorra vonatkozó explicit megoldás diszkrét formája, amit a szimuláció során alkalmaztunk:

[latexmath]
++++
S(t_k) = S(0) \exp\left( \left(\mu - \frac{\sigma^2}{2}\right)t_k + \sigma W(t_k) \right)
++++

ahol stem:[W(t_k)] a diszkrét Brown-mozgás (véletlen bolyongás) értéke a stem:[t_k] időpontban.

==== Adatok előkészítése

A modellhez a 2001 és 2024 közötti, első kerületre vonatkozó napi bűnesetszámokat tartalmazó historikus adatokat használtuk. A 2025-ös év adatait validációs célokra különítettük el.

==== Paraméterek definiálása és becslése

A szimulációhoz szükséges paramétereket a historikus adatokból becsültük, a notebookban leírtak szerint:

*   **stem:[S_0] (Kezdeti érték):** A historikus idősor utolsó ismert értéke, azaz a 2024-es év utolsó napi bűnesetszáma.
*   **stem:[\Delta t] (Időlépés):** Az adatok gyakoriságának megfelelő időlépés. Esetünkben havi adatokkal dolgoztunk, így stem:[\Delta t = 1] (hónap).
*   **stem:[T] (Előrejelzési horizont):** Az előrejelzés teljes időtávja. A 2025-ös év első 9 hónapjára készítettünk becslést, így stem:[T = 9] hónap.
*   **stem:[N] (Időpontok száma):** Az előrejelzési horizonton belüli időpontok száma, stem:[N = T / \Delta t = 9].
*   **stem:[\mu] (Drift paraméter):** A historikus adatok relatív változásainak átlaga. Ez a paraméter a bűnesetek számának átlagos növekedési vagy csökkenési rátáját méri.

[latexmath]
++++
\mu = \text{mean}\left(\frac{S_i - S_{i-1}}{S_{i-1}}\right)
++++
*   **stem:[\sigma] (Volatilitás):** A historikus adatok relatív változásainak szórása. Ez a paraméter az adatsor ingadozásának mértékét, a véletlen sokkok nagyságát határozza meg.

[latexmath]
++++
\sigma = \text{std}\left(\frac{S_i - S_{i-1}}{S_{i-1}}\right)
++++
*   **stem:[b] (Brown-növekmények):** Egy vektor, amely minden egyes időlépéshez egy standard normális eloszlásból stem:[({N}(0,1))] húzott véletlen számot tárol. Ezek adják a modell sztochasztikus viselkedésének alapját.
*   **stem:[W] (Brown-pálya):** A Brown-növekmények kumulatív összege. A stem:[W(t_k)] érték a teljes véletlen hatást összegzi a szimuláció kezdetétől a stem:[t_k] időpontig.

==== Szimuláció és eredmények kiértékelése

A szimuláció során több lehetséges jövőbeli pályát (szcenáriót) generáltunk, hogy felmérjük a modell bizonytalanságát. Minden szcenárióhoz egyedi Brown-pályát (stem:[W]) generáltunk.

Az előrejelzett értékeket a fentebb bemutatott diszkrét megoldási képlettel számítottuk ki minden egyes szcenárióra és időpontra.

Az eredmények kiértékeléséhez az előrejelzett értékeket összevetettük a 2025-ös év tényleges (validációs) adataival. A modell pontosságát minden szcenárióra a gyökös négyzetes hiba (RMSE) metrikával mértük:

[latexmath]
++++
\mathrm{RMSE} = \sqrt{\frac{1}{N} \sum_{i=1}^N (y_i - \hat{y}_i)^2}
++++

ahol stem:[y_i] a tényleges, stem:[\hat{y}_i] pedig az előrejelzett érték.

A notebookban az elemzést kiterjesztettük az összes kerületre. Minden kerületre külön-külön futtattuk a modellt, és kiválasztottuk a legalacsonyabb RMSE értékkel rendelkező szcenáriót, mint a legvalószínűbb előrejelzést. A vizualizációk és a hibamértékek alapján a geometriai Brown-mozgás modell alkalmasnak bizonyult a bűnesetek számának sztochasztikus modellezésére és előrejelzésére.


\end{document}
